\documentclass[a4paper,12pt,twoside,openright,titlepage]{article}

%Additional packages
\input{styles/packages}

% Listing style
\input{styles/lstset}

% Uncomment for production
% \syntaxonly

% Style
\pagestyle{headings}

% Define document
\author{D. Leeuw}
\title{Python System Checker}
\date{\today\\
1.0.0 \\
\vfill
\raggedright
\copyright\ 2026 Dennis Leeuw\\
\input{styles/licentie-titlepage}}


\begin{document}
\selectlanguage{dutch}

\maketitle

%%%%%%%%%%%%%%%%%%%
%%% Introductie %%%
%%%%%%%%%%%%%%%%%%%

%B% \frontmatter
\section{Over dit Document}
\subsection{Leerdoelen}
\input{src/opdracht-functies-system_status-leerdoelen}
\subsection{Voorkennis}
\input{src/opdracht-functies-system_status-voorkennis}

%%%%%%%%%%%%%%%%%
%%% De inhoud %%%
%%%%%%%%%%%%%%%%%
%%% \tableofcontents

%B% \mainmatter

\section{Situatie}
\input{src/opdracht-functies-system_status-situatie}

\section{Opdracht: Systeem check}

\subsection{Maak controle functies}
\input{src/opdracht-functies-system_status-opdracht}
\subsubsection{Te maken functies}
\input{src/opdracht-functies-system_status-opdracht_functies}

\subsection{Waarschuwen bij problemen}
\input{src/opdracht-functies-system_status-waarschuwing}
\subsubsection{De status functie}
\input{src/opdracht-functies-system_status-controle_functie}

\subsection{Systeem status}
\input{src/opdracht-functies-system_status-totaal}

\subsection{Het hoofdprogramma}
\input{src/opdracht-functies-system_status-main}

\subsection{Maak een menu}
\input{src/opdracht-functies-system_status-menu}

\subsection{Vragen}
\input{src/opdracht-functies-system_status-vragen}

%%%%%%%%%%%%%%%%%%%%%
%%% Index and End %%%
%%%%%%%%%%%%%%%%%%%%%
%\backmatter
\printindex
\end{document}

%%% Last line %%%
